\section{Six-role lower bounds}\label{sec:lower-bounds}



The exact-six question of
\S\ref{subsec:scoped-exact-six-question} has two directions. The
lower-bound direction asserts that each of \Pone{}, \ldots, \Psix{}
is instantiated in \FATCD{}; the upper-bound classification of
Section~\ref{sec:upper-bound-classification} will assert, under a recognition
hypothesis, that no seventh role is needed in the covered scope. This section establishes the
lower-bound direction by constructing one admissible witness for each
role. The content is occurrence, not exhaustion: lower bounds alone
do not answer the exact-six question, but they rule out the
trivialization in which some role has no admissible witness in the
domain.

\subsection{Lower-bound theorem and common witness requirements}\label{subsec:lower-bound-theorem}

\begin{theorem}[Six-role lower bounds]\label{thm:six-role-lower-bounds}
For each role index \(i \in \{1,\ldots,6\}\), there exists a description
\(D_i \in \FATCD{}\) and a derived active projection
\(\pi_i(D_i)\) witnessing the role channel \(P_i\) in the sense of
Definition~\ref{def:role-projection}. In particular, each of
\Pone{}, \Ptwo{}, \Pthree{}, \Pfour{}, \Pfive{}, \Psix{} is instantiated
by at least one admissible description in \FATCD{}. This is an occurrence
theorem, not an exhaustion theorem.
\end{theorem}

A witness for role \(P_i\) consists of the data specified below,
following the role-projection discipline of
\S\ref{subsec:derived-role-projections} and the threshold-attachment
discipline of \S\ref{subsec:threshold-attachment}.

\begin{definition}[Common witness data for role \(P_i\)]\label{def:common-witness-data}
A \emph{witness of role \(P_i\) in \FATCD{}} is a tuple
\[
\bigl(D_i,\; C_i,\; \Theta_i,\; W_i,\; \mathrm{Level}_i,\;
v_i^{\mathrm{fin}},\; v_i^{\mathrm{cl}},\; v_i^{\mathrm{aud}},\;
n_i,\; \pi_i(D_i)\bigr)
\]
where
\begin{itemize}[topsep=2pt,itemsep=1pt]
  \item \(D_i \in \FATCD{}\) is the ambient admissible description;
  \item \(C_i \subseteq \Raw{D_i}\) is the raw feature cluster on which
    the projection is interpreted;
  \item \(\Theta_i\) is the threshold annotation attached to \(C_i\);
  \item \(W_i\) is the witness schema for the role meaning of \(P_i\)
    fixed in the \emph{Witness types} table of
    Section~\ref{sec:primitive-role-architecture};
  \item \(\mathrm{Level}_i\) is the level assignment in the sense of
    Definition~\ref{def:level-trichotomy};
  \item \(v_i^{\mathrm{fin}}, v_i^{\mathrm{cl}}, v_i^{\mathrm{aud}}\) are
    the statements that \(D_i\) satisfies Definitions~\ref{def:finite},
    \ref{def:closed} and~\ref{def:audited} respectively, each with its proof;
  \item \(n_i\) is the explicit nonclaim record carried by the witness;
  \item \(\pi_i(D_i)\) is the derived role projection of
    Definition~\ref{def:role-projection}.
\end{itemize}
The witness is \emph{admissible} when \(v_i^{\mathrm{fin}}\),
\(v_i^{\mathrm{cl}}\) and \(v_i^{\mathrm{aud}}\) show \(D_i\in\FATCD{}\) and
\(\pi_i(D_i,C_i)\) is admissible in the sense of
Definition~\ref{def:role-projection}, with \(\Theta_i\), the \(W_i\)
witness-schema annotation, \(\mathrm{Level}_i\) and \(n_i\) as its threshold,
witness-schema, level and nonclaim annotations; in particular \(W_i(C_i)\)
holds.
\end{definition}

The proof of Theorem~\ref{thm:six-role-lower-bounds} proceeds role by
role: for each \(i \in \{1,\ldots,6\}\) we exhibit an admissible
tuple of the form in Definition~\ref{def:common-witness-data}. The
six constructions are collected in
\S\ref{subsec:per-role-witnesses}.

\subsection{Per-role witnesses}\label{subsec:per-role-witnesses}

Each construction below specifies a raw feature cluster, threshold
annotation, witness schema, level assignment, derived active
projection, and explicit nonclaim record. The finite-description,
closure, and audit/bookkeeping verifications are routine once these
data are placed in a description built from the raw grammar of
\S\ref{subsec:raw-grammar} and checked against the admissibility
conditions of \S\ref{subsec:conditions}. In the \Pone{} construction,
\(\sim\) is a relation on the source records and \(\sim'\) a relation on
the target records; descent asks whether \(x\sim y\) implies
\(F(x)\sim' F(y)\). In the constructions, \(\mathrm{Ann}_\Theta[t]\) is a
threshold annotation whose recorded label is \(t\), and
\(v_i^{\mathrm{fin}}, v_i^{\mathrm{cl}}, v_i^{\mathrm{aud}}\) are the checks
that \(D_i\) satisfies Definitions~\ref{def:finite}, \ref{def:closed}
and~\ref{def:audited}. \par\emph{Annotations.} In each construction, \(D_i\) includes a claim
annotation naming \(C_i\), \(W_i\), the asserted check result, hypotheses and
scope, and an audit annotation referencing it and the finite evidence used by
the check, together with the threshold, level, witness-schema, collapse-case
and nonclaim annotations as declared records; every endpoint or reference used in a
construction is included as a declared record of that \(D_i\), even where the
abbreviated display names only its principal records. In each construction,
the claim explicitly names \(P_i\) and \(C_i\), and the threshold,
witness-schema, collapse-case, level and nonclaim annotations each have a
target field listing exactly \(C_i\). The verification \(v_i^{\mathrm{aud}}\) is the check that
these entries are honest in the sense of Definition~\ref{def:audited}, not
merely present. \par\emph{Fields.} Each \(D_i\) contains only the displayed witness records, their declared
inputs, and the stated annotations. Every content record either contributes a
required field to an aligned subcluster of \(C_i\), carries a covered \(P_i\)
seed, or is role-neutral data without an assertion of its own, such as a record named
in a declared domain, codomain, source or target field of a required-field
contributor, whether or not a check reads it. No content record
makes an assertion independent of the displayed witness. Relation records
other than \(\mathrm{Real}\) in \(D_2\), and map, path, index and state
records, carry only datum and reference fields; every law,
invariant or comparison is a field of the displayed check, comparison or
law-bearing record. Each threshold annotation \(\mathrm{Ann}_\Theta[t]\) carries
only its typing, target and label \(t\), and each level annotation carries only
its typing, target and level; each claim annotation carries, besides its
statement and hypothesis fields, only datum and reference fields; and each
collapse-case, nonclaim, promotion and translation annotation carries only
datum and reference fields. \par\emph{Audits.} In each \(D_i\), the audit of the \(W_i\) claim
lists the row's content certificates carried by members of \(C_i\), each
recomputed true, and a certificate \((W_i(C_i),\mathrm{true})\); its conclusion
is \(W_i(C_i)\). The latter certificate is recomputed after the content checks
and, for \(i=6\), after \(a_\tau\) is established honest in
\(\mathrm{Audit}_0(D_6)\). The conclusion thus follows from a signed premise,
and the audit is honest.

\begin{proposition}[\Pone{} lower bound]\label{prop:p1-witness}
There exists \(D_1 \in \FATCD{}\) with a derived active projection
\(\pi_1(D_1)\) witnessing \Pone{}.
\end{proposition}

\begin{proof}
Let \(D_1\) carry the finite state set \(S=\{x,y,x',y'\}\) with four distinct
state records and a declared object record \(S\) listing \(x,y,x',y'\), the source equivalence relation listing the classes \(\{x,y\}\), \(\{x'\}\), \(\{y'\}\), and the target equality relation on \(S\), listing the four singleton classes, both relation records
of declared use quotient. Let \(F:S\to S\) be a
map record whose declared use is \emph{update}, with \(F(x)=x'\),
\(F(y)=y'\), \(F(x')=x'\), and \(F(y')=y'\). A comparison record checks
\(x\sim y\) but \(F(x)\not\sim'F(y)\); hence the related source pair witnesses
the descent obstruction. Attach a threshold annotation
\(\Theta_1 = \mathrm{Ann}_\Theta[\text{descent obstruction}]\) to the
cluster \(C_1\) comprising these records, and record a witness-schema annotation naming \Pone{} and the \(W_1\) row;
\(W_1(C_1)\) holds by alternative (b), since the comparison record names
\((x,y)\) and carries the certificate
\((x\sim y\wedge((x\in\mathrm{dom}F\leftrightarrow y\notin\mathrm{dom}F)\vee(x\in\mathrm{dom}F\wedge y\in\mathrm{dom}F\wedge F(x)\sim'F(x)\wedge F(y)\sim'F(y)\wedge\neg F(x)\sim'F(y))),\mathrm{true})\),
whose second disjunct recomputes true. Take \(\mathrm{Level}_1\) on the verification
level, and let \(\pi_1(D_1)\) be the derived active projection
interpreting \(C_1\) via \((\Theta_1,W_1,\mathrm{Level}_1)\). The
finiteness, closure, and audit verifications \(v_1^{\mathrm{fin}},
v_1^{\mathrm{cl}}, v_1^{\mathrm{aud}}\) then attach routinely: the claim has
statement \(W_1(C_1)\); its audit references the claim and every member of \(C_1\), lists the comparison certificate and
\((W_1(C_1),\mathrm{true})\), both recomputed true, and concludes \(W_1(C_1)\).
The descriptions \(D_2\)--\(D_6\) follow the same template; for \(D_6\), the
event-or-trace audit is established honest in \(\mathrm{Audit}_0(D_6)\) before
\((W_6(C_6),\mathrm{true})\) is recomputed. The
nonclaim record \(n_1\) excludes any \Pfive{}-packaging
identification (Remark~\ref{rem:repaired-alignments}) and any claim of
an unconditional descended map in the obstruction case.
\end{proof}

In the running example, \(\sim\) and \(\sim'\) are both the grouping by
\(q\), \(x=0\) and \(y=1\): \(q(0)=q(1)\), while \(q(F0)=B\neq A=q(F1)\).

\begin{proposition}[\Ptwo{} lower bound]\label{prop:p2-witness}
There exists \(D_2 \in \FATCD{}\) with a derived active projection
\(\pi_2(D_2)\) witnessing \Ptwo{}.
\end{proposition}

\begin{proof}
Let \(D_2\) carry the state records \(0,1,2,3\), the candidate class of
realizers; an object record \(S_2\) listing \(0,1,2,3\), declared as the domain
of \(q\) and as the candidate class of \(\mathrm{Real}\); a quotient object record \(\{A,B\}\) and a map record \(q\) of
declared use packaging with \(q(0)=q(1)=A\) and \(q(2)=q(3)=B\), whose kernel
\(E\) the comparison uses; a predicate record \(s\) declared as a
macro-specification with formula \(s(r)\equiv(r=0\lor r=2)\); a relation record
\(\mathrm{Real}=\{(0,s),(2,s)\}\), of declared use \emph{realization}, with
candidate class \(K=\{0,1,2,3\}\) and target \(s\),
whose datum fields record the formula's values true, false, true, false at \(0,1,2,3\) and which carries
the certificate \((s(0),\mathrm{true})\) exhibiting the satisfying realizer
\(0\); and a comparison record carrying the certificate
\((q(0)=q(1)\wedge s(0)\wedge\neg s(1),\mathrm{true})\), certifying that the
\(E\)-class \(\{0,1\}\) meets both the realizer set \(\{0,2\}\) and its
complement \(\{1,3\}\) in \(K\), since \(0\in\{0,2\}\) and \(1\notin\{0,2\}\)
although \(q(0)=q(1)\), so that
realizability is not reducible to quotient-class equality alone; the finite quotient-class comparison data are
fields of this comparison record, and any separately declared class
identifiers belong to \(C_2\). Attach
\(\Theta_2 = \mathrm{Ann}_\Theta[\text{nontrivial representability}]\)
to the cluster \(C_2\) of the records \(0,1,2,3\), \(S_2\), \(q\), \(\{A,B\}\), \(s\), \(\mathrm{Real}\) and the comparison record (the comparison reads no class-identifier record; the object records \(A\), \(B\), which occur only as values of \(q\) and are listed by \(\{A,B\}\), are declared records of \(D_2\) outside \(C_2\), have no assertion field and satisfy (R3)), and record a witness-schema annotation naming \Ptwo{} and the \(W_2\) row,
and \(W_2(C_2)\) holds by the satisfying branch: \(s\) has the single free
variable \(r\), \(\mathrm{Real}\) records \(s\) at each member of \(K\) and
exhibits the realizer \(0\), and the comparison certifies that the class
\(\{0,1\}\) of \(\ker q\), whose domain \(S_2\) contains \(K\), meets both
\(\{0,2\}\) and \(\{1,3\}\). Take
\(\mathrm{Level}_2\) on the verification level, and let \(\pi_2(D_2)\)
be the derived active projection determined by
\((C_2,\Theta_2,W_2,\mathrm{Level}_2)\). The finiteness, closure, and
audit checks then follow from the construction. The
nonclaim \(n_2\) excludes collapse to \Pfive{} quotient equality alone
and excludes arbitrary relation promotion.
\end{proof}

The states and the map \(q\) are those of the running example.

\begin{proposition}[\Pthree{} lower bound]\label{prop:p3-witness}
There exists \(D_3 \in \FATCD{}\) with a derived active projection
\(\pi_3(D_3)\) witnessing \Pthree{}.
\end{proposition}

\begin{proof}
Let \(D_3\) carry a state record \(u\) and an object record \(U\) listing \(u\);
generating moves \(a,b\), recorded as map records of declared use
\emph{other} with declared source and target \(U\) and \(a(u)=b(u)=u\);
two path records \(\rho_1=ab\) and \(\rho_2=ba\), each listing the successive
records \((u,u,u)\), both from \(u\) to \(u\) and each referencing the
comparison record; and a comparison record referencing both \(\rho_1\) and \(\rho_2\) and carrying the
rewrite rule \(aa\to a\) as a rule field (Definition~\ref{def:raw-record}) holding the pair of words \((aa,a)\), the finite monoid \(N=\{e,\bar a,\bar b\}\), its elements recorded as the rational constants \(e=0\), \(\bar a=1\), \(\bar b=2\) in a list field and its multiplication as a finite table field, with
\(ey=y\) and \(xy=x\) for \(x\neq e\), the images \(\varphi(a)=\bar a\) and
\(\varphi(b)=\bar b\), a move-list datum field \(M=(a,b)\), and a certificate whose
formula is the full prescribed \(W_3\)(a) condition, reading \(R,N,e,m,\varphi,M\)
from the comparison record and \(w_1,w_2\) from the route records. The
assignments \(R=((aa,a))\), \(M=(a,b)\), \(w_1=ab\), \(w_2=ba\) specify their
recorded values, not literal substitutions in the certificate. It checks that \(e\in N\), that \(N\) is closed under \(m\) and
associative with identity \(e\), that \(\varphi(a),\varphi(b)\in N\), that \(\varphi(\lambda)=\varphi(\mu)\) for each pair \((\lambda,\mu)\) listed in that datum field (here \(\varphi(aa)=\bar a=\varphi(a)\)), and that
\(\varphi^{m,e}(w_{\rho_1})\ne\varphi^{m,e}(w_{\rho_2})\), where \(w_{\rho_j}\) is the recorded word of \(\rho_j\) (here \(ab\) and \(ba\)) and \(m\) is
the recorded multiplication table. On the displayed table,
\(\varphi^{m,e}(ab)=\bar a\) and
\(\varphi^{m,e}(ba)=\bar b\). Let
\(C_3=\{u,U,a,b,\rho_1,\rho_2,\text{comparison}\}\).
Attach \(\Theta_3 = \mathrm{Ann}_\Theta[\text{route mismatch}]\) to the
cluster \(C_3\), and
record a witness-schema annotation naming \Pthree{} and the \(W_3\) row;
\(W_3(C_3)\) holds because the two generated same-source, same-target route
records are distinguished by a finite-monoid non-reducibility
certificate of the \(W_3\) row. Take \(\mathrm{Level}_3\) on the
verification level, and let \(\pi_3(D_3)\) be the derived active
projection determined by
\((C_3,\Theta_3,W_3,\mathrm{Level}_3)\). The routine verifications
\(v_3^{\mathrm{fin}}, v_3^{\mathrm{cl}}, v_3^{\mathrm{aud}}\) then
follow. The nonclaim \(n_3\) excludes conflation with
\Pone{} route undefinedness and with direct cross-level comparison.
Routes are compared in the category presented by the recorded moves and rules,
not by equality of composite maps: here \(a\) and \(b\) have the same graph on
\(\{u\}\), and the witness certifies only that the recorded rules do not
identify \(ab\) with \(ba\).
\end{proof}

The monoid is the first-move monoid of
Section~\ref{sec:primitive-role-architecture}.

\begin{proposition}[\Pfour{} lower bound]\label{prop:p4-witness}
There exists \(D_4 \in \FATCD{}\) with a derived active projection
\(\pi_4(D_4)\) witnessing \Pfour{}.
\end{proposition}

\begin{proof}
Let \(D_4\) carry index records \(0,\tfrac12,1,2\); object records
\(\{0\},\{1\},\{2\}\) listing one index record each, and object records
\(I=\{0,1,2\}\) and \(I'=\{0,\tfrac12,1,2\}\); order records listing the pairs \((0,1),(0,2),(1,2)\)
on \(I\) and all six pairs of \(0<\tfrac12<1<2\) on \(I'\); transition records
\(t_{01}:\{0\}\to\{1\}\) with \(t_{01}(0)=1\) and \(t_{12}:\{1\}\to\{2\}\) with
\(t_{12}(1)=2\); and a refinement record
for the inclusion \(I\hookrightarrow I'\).
Every endpoint and index reference resolves to a declared record. Each
order, transition and refinement record carries its certificate as a law
field: finite certificates check that both orders are irreflexive and
transitive and each relates at least one pair, that each transition is defined on its declared domain with values
in its declared codomain, and that the refinement maps \(I\) into \(I'\) and preserves order, each certificate being the prescribed law formula of Section~\ref{sec:primitive-role-architecture} for its record. Attach
\(\Theta_4=\mathrm{Ann}_\Theta[\text{nontrivial stage/refinement}]\)
to the cluster \(C_4\) comprising these records, and record a
witness-schema annotation naming \Pfour{} and the \(W_4\) row; \(W_4(C_4)\)
holds because the records carry declared endpoints or indices and a law for
them. Take
\(\mathrm{Level}_4\) on the structural-categorical level, and let
\(\pi_4(D_4)\) be the derived active projection determined by
\((C_4,\Theta_4,W_4,\mathrm{Level}_4)\). The routine verifications
\(v_4^{\mathrm{fin}}, v_4^{\mathrm{cl}}, v_4^{\mathrm{aud}}\) then
follow. The nonclaim \(n_4\) excludes bare labels and arbitrary
relations or transitions without laws.
\end{proof}

In the running example the transition record for \(F\) is present, but no
threshold annotation is attached to it, so the \Pfour{} channel of
\(D_{\mathrm{ex}}\) is \(\belowthr\) (Section~\ref{sec:decomposition-exhaustion}).
Only the attachment data are missing there: the cluster formed by the transition
record of \(F\), \(S_0\) and the four state records already satisfies \(W_4\). The description \(D_4\) carries a threshold, a
witness-schema and a level annotation on a larger law-bearing cluster, a stage
tower \(0<1<2\) with transitions and a refinement law.

\begin{proposition}[\Pfive{} lower bound]\label{prop:p5-witness}
There exists \(D_5 \in \FATCD{}\) with a derived active projection
\(\pi_5(D_5)\) witnessing \Pfive{}.
\end{proposition}

\begin{proof}
Take two distinct declared object records \(x,y\) (the micro-elements), an object record
listing them as the declared domain of \(q\), a declared state record \(*\) and a
quotient object record \(Q\) listing it, the relation record \(E\), of declared use
quotient, relating both, and the map record \(q:\{x,y\}\to Q\), with declared use \emph{packaging}, sending both to \(*\). Writing \(X=\{x,y\}\) for the declared domain, the finite check
record carries the certificate, recomputed true,
\begin{multline*}\forall u,v\in X\,(E(u,v)\leftrightarrow q(u)=q(v))\wedge x\ne y\wedge q(x)=q(y)\\ \wedge\forall z\in Q\,\exists u\in X\,(q(u)=z)\wedge\forall u\in|E|\,(u\in X)\wedge\forall u\in X\,(q(u)\in Q).\end{multline*} Thus
\(x\ne y\), \(q(x)=q(y)\), and the complete quotient-validity condition in the
\(W_5\) row holds. Let \(C_5\) comprise these declared records and this check,
and let \(D_5\) carry them. Attach
\(\Theta_5 = \mathrm{Ann}_\Theta[\text{nontrivial quotient packaging}]\)
to the cluster \(C_5\) comprising these records, and record a
witness-schema annotation naming \Pfive{} and the \(W_5\) row; \(W_5(C_5)\)
holds because distinct raw objects are related and packaged together by a
packaging map with quotient-validity evidence. Take
\(\mathrm{Level}_5\) on the verification level, and let \(\pi_5(D_5)\)
be the derived active projection determined by
\((C_5,\Theta_5,W_5,\mathrm{Level}_5)\). The routine verifications
\(v_5^{\mathrm{fin}}, v_5^{\mathrm{cl}}, v_5^{\mathrm{aud}}\) then
follow. The nonclaim \(n_5\) excludes identity packaging and
witnesses whose identified records are not distinct.
\end{proof}

\begin{proposition}[\Psix{} lower bound]\label{prop:p6-witness}
There exists \(D_6 \in \FATCD{}\) with a derived active projection
\(\pi_6(D_6)\) witnessing \Psix{}.
\end{proposition}

\begin{proof}
Let \(D_6\) consist of state records \(s_0,s_1,s_2\); an event record \(e\)
referencing \(s_0\); a trace record \(\tau\) listing \((s_0,s_1,s_2)\) and
referencing \(e\); a provenance record referencing \(e\) and \(\tau\); a map
record \(t:\{s_0,s_1\}\to\{s_0,s_1,s_2\}\) of declared use other, with declared
object records for its source and target, and \(t(s_0)=s_1\), \(t(s_1)=s_2\); and a
check record referencing \(\tau\) and \(t\) that carries the certificate
\(\tau_1=t(\tau_0)\wedge\tau_2=t(\tau_1)\), recomputed true; together with the annotations listed below. Let \(C_6\) consist of \(s_0,s_1,s_2\), \(e\), \(\tau\), the provenance record, \(t\), the source and target object records of \(t\), and the check record. Record a claim annotation \(c_\tau\) referencing \(e\), \(\tau\) and the check record, whose statement is the formula \(\tau_1=t(\tau_0)\wedge\tau_2=t(\tau_1)\) of the check record's certificate (so \(c_\tau\) satisfies the representation condition of Definition~\ref{def:annotations}), with an audit \(a_\tau\) referencing
\(c_\tau\), \(e\), \(\tau\), the provenance record, \(t\) and the check record and certifying
precisely those equations from the check record; both are assigned to the
active \Psix{} projection; then record the claim naming \Psix{} and
\(C_6\) with a second audit \(a_6\) certifying \(W_6(C_6)\), whose
\(\Audit{D}\) conjunct is discharged by \(a_\tau\). Attach
\(\Theta_6 = \mathrm{Ann}_\Theta[\text{nontrivial audit schema}]\) to
the cluster \(C_6\), and record a
witness-schema annotation naming \Psix{} and the \(W_6\) row; \(W_6(C_6)\)
holds because \(C_6\) contains the event \(e\), the trace \(\tau\), which
references \(e\), and the check record, which references \(\tau\) and carries
\(\tau_1=t(\tau_0)\wedge\tau_2=t(\tau_1)\), exactly the conjunction of
successive-entry replay equalities that the \(W_6\) row requires; and
\(a_\tau\in\mathrm{Audit}_0(D_6)\) refers to \(c_\tau\), \(e\), \(\tau\) and the check
record and concludes that formula. Take \(\mathrm{Level}_6\) on the verification
level, and let \(\pi_6(D_6)\) be the derived active projection
determined by \((C_6,\Theta_6,W_6,\mathrm{Level}_6)\). The routine
verifications \(v_6^{\mathrm{fin}}, v_6^{\mathrm{cl}},
v_6^{\mathrm{aud}}\) then follow. The nonclaim \(n_6\) excludes passive
logs, certificate-only content, and partial audit records.
\end{proof}

Propositions~\ref{prop:p1-witness}--\ref{prop:p6-witness} supply, for
each \(i \in \{1,\ldots,6\}\), an admissible witness tuple of the
form in Definition~\ref{def:common-witness-data}. This proves
Theorem~\ref{thm:six-role-lower-bounds}.

The running example \(D_{\mathrm{ex}}\) of Section~\ref{sec:fatcd-domain}
is an explicit description of this kind for three roles: its split pair
witnesses \Pone{}, its quotient \(\{A,B\}\) with packaging map \(q\)
witnesses \Pfive{}, and its trace with replay and check witnesses \Psix{}.

\begin{remark}[What the lower bounds do and do not establish]\label{rem:lower-bounds-scope}
Theorem~\ref{thm:six-role-lower-bounds} asserts that each of the six
roles is instantiated somewhere in \FATCD{}. It does not prove
exact-six, does not rule out candidate seventh-role threats, does not
prove decomposition, and does not prove uniqueness of any later
decomposition/exhaustion record. It also does not assert that any
single description activates all six channels simultaneously;
Proposition~\ref{prop:channel-vs-activation} (case \(k=6\)) constructs one. Those
questions are handled in
Sections~\ref{sec:upper-bound-classification}
and~\ref{sec:decomposition-exhaustion}.
\end{remark}
